\chapter*{General Introduction}
\addcontentsline{toc}{chapter}{General Introduction}
\markboth{General Introduction}{}
\setcounter{page}{1}

% STRUCTURE: Context -> Problem -> Approach -> Outline (one paragraph each, 4 to 8 lines).
% TO BE WRITTEN LAST, once the chapters have been finalized.

% 1. CONTEXT: firewalls, traffic logs, data volume, and the value of automated analysis.
\TODO{Context (1 paragraph): the role of a firewall and its logs, and the value of analyzing them using data mining}

% 2. PROBLEM: predicting the recorded action (allow, deny, drop, reset-both) from flow attributes.
%    Clarify that this is a retrospective log classification task, not a protection system.
\TODO{Problem Statement (1 paragraph): the research question and challenges (class imbalance, duplicates, data leakage)}

% 3. APPROACH: UCI dataset, methodology, six conventional classifiers, rigorous evaluation, and Streamlit application.
\TODO{Approach and Contributions (1 paragraph)}

% 4. OUTLINE: one sentence per chapter, using \ref (chapter numbers update automatically).
To facilitate reading, this report is organized into six chapters:
\begin{itemize}
  \item[$-$] \textbf{Chapter~\ref{chap:contexte}} presents the context and problem statement.
  \item[$-$] \textbf{Chapter~\ref{chap:fondements}} reviews the theoretical foundations.
  \item[$-$] \textbf{Chapter~\ref{chap:donnees}} describes the dataset and its exploratory analysis.
  \item[$-$] \textbf{Chapter~\ref{chap:methodologie}} details the experimental methodology.
  \item[$-$] \textbf{Chapter~\ref{chap:resultats}} presents and discusses the results.
  \item[$-$] \textbf{Chapter~\ref{chap:application}} describes the interactive application.
\end{itemize}